% 交直流混合最优潮流：Parity IPM 求解核
% 本章文件由 \input 引入，不含导言区。
\section{交直流混合最优潮流：Parity IPM 求解核}\label{sec:parityipm}

本章给出交直流混合 OPF（第~\ref{sec:parityform}~章建模）的自研原始对偶内点
求解核（下称 Parity IPM 核）的数值口径：它以 Parity 建模层装配出的
\srcpath{parity::Problem} 为输入，在同一个牛顿框架内同时推进原始变量、
松弛变量与全部乘子，输出收敛证书、状态向量与诊断计数。全章符号锚点均指
\srcpath{src/optimal\_power\_flow/parity\_ipm.cpp} 与
\srcpath{include/hacdcpf/optimal\_power\_flow/native\_ipm\_solver.hpp}
（IPMOptions/IPMResult 类型定义），驱动层锚点指
\srcpath{src/optimal\_power\_flow/ac\_opf.cpp}。核内全部方程为标幺值
（$S_{\mathrm{base}}=100$~MVA，见 \ref{sec:parityform}~节）；成本币种未定。

\subsection{求解核定位与调用链}

\subsubsection*{功能定位}
Parity IPM 核是平台 AC OPF 三条非线性求解路线之一（另两条为
\ref{sec:acopfnative}~节的 Native AC 外部障碍环路与本章
\ref{subsec:pipm-ipopt}~节的嵌入式 Ipopt 适配）。与 Native AC 路径的
“外层障碍参数 $+$ 内层牛顿”两层结构不同，本核在 Parity 建模层给出的
全空间问题上直接运行 Mehrotra 预测--校正原始对偶内点法：所有变量块与
约束行在 \srcpath{Problem::vidx}/\srcpath{Problem::cidx} 中显式编号
（\srcpath{include/hacdcpf/optimal\_power\_flow/formulation.hpp:VarIndex、:ConstraintIndex}），
核不再做任何设备级聚合或惩罚逼近，因此它是混合 AC/DC 统一建模的
首选审查路径（\srcpath{include/hacdcpf/optimal\_power\_flow/parity\_ipm\_solver.hpp:solve\_ac\_opf\_parity}
头文件注释）。核的公开入口为
\srcpath{solve\_primal\_dual\_ipm}（声明
\srcpath{native\_ipm\_solver.hpp:solve\_primal\_dual\_ipm}，定义 :solve\_primal\_dual\_ipm），配套的同伦切线
例程 \srcpath{homotopy\_tangent} 复用同一套 KKT 机器为 Davidenko 目标
同伦提供路径切向量（声明 \srcpath{native\_ipm\_solver.hpp:homotopy\_tangent}，
定义 :homotopy\_tangent）。

\subsubsection*{调用链与驱动层}
头文件 \srcpath{parity\_ipm\_solver.hpp:solve\_ac\_opf\_parity} 声明了高层驱动
\srcpath{solve\_ac\_opf\_parity}，但该声明在本检出处\textbf{没有对应定义}
（全仓检索仅命中声明本身），属悬置声明；实际驱动是
\srcpath{ac\_opf.cpp} 中的文件静态函数 \srcpath{solve\_with\_parity\_ipm}
（src/optimal\_power\_flow/ac\_opf.cpp:solve\_with\_parity\_ipm），由公共门面 \srcpath{solve\_ac\_opf}
（src/optimal\_power\_flow/ac\_opf.cpp:solve\_ac\_opf）按 \fld{ACOPFOptions::ac\_solver\_backend} 分派：
\begin{itemize}
  \item \fld{ACOPFSolverBackend::ParityIPM}：强制本核
        （内层枚举 \fld{ParityInnerSolver::NativeIPM}，src/optimal\_power\_flow/ac\_opf.cpp:solve\_ac\_opf\_impl）；
  \item \fld{ACOPFSolverBackend::Ipopt}：强制嵌入式 Ipopt
        （src/optimal\_power\_flow/ac\_opf.cpp:solve\_ac\_opf\_impl）；
  \item \fld{Auto}（默认）：本核先行，未收敛且 \fld{allow\_fallback}
        为真时回退 Ipopt（src/optimal\_power\_flow/ac\_opf.cpp:solve\_with\_parity\_ipm）；凡检测到 DC/VSC/LCC
        子系统即自动启用本核所在的 Parity 路径（src/optimal\_power\_flow/ac\_opf.cpp:solve\_ac\_opf\_impl）。
\end{itemize}
内层求解器三态枚举 \fld{ParityInnerSolver} 定义于 src/optimal\_power\_flow/ac\_opf.cpp:ParityInnerSolver。
驱动层还负责把面向用户的 \fld{ACOPFOptions} 折算为核参数
\fld{IPMOptions}（src/optimal\_power\_flow/ac\_opf.cpp:solve\_with\_parity\_ipm，映射关系见
\ref{subsec:pipm-params}~节参数表说明列），并在求解后把 \fld{IPMResult}
归因回 \fld{ACOPFResult}（src/optimal\_power\_flow/ac\_opf.cpp:solve\_with\_parity\_ipm，见
\ref{subsec:pipm-errors}~节）。整体调用链与门控位置见图~\ref{fig:pipm-chain}。

\begin{figure}[H]
\centering
\resizebox{\textwidth}{!}{%
\begin{tikzpicture}[
  font=\small,
  box/.style={draw, rectangle, align=center, minimum height=0.85cm,
              inner sep=4pt},
  lab/.style={font=\footnotesize, align=center},
  arr/.style={-{Stealth[length=2.2mm]}, thick},
  node distance=0.55cm and 0.9cm]
  \node[box] (api) {\srcpath{solve\_ac\_opf}\\[-2pt]{\footnotesize src/optimal\_power\_flow/ac\_opf.cpp:solve\_ac\_opf}};
  \node[box, below=of api] (drv) {\srcpath{solve\_with\_parity\_ipm}\\[-2pt]
        {\footnotesize src/optimal\_power\_flow/ac\_opf.cpp:solve\_with\_parity\_ipm（后端分派/选项折算/结果归因）}};
  \node[box, below left=0.7cm and 0.2cm of drv] (prob) {Parity 建模层\\
        \srcpath{build\_problem}\\[-2pt]{\footnotesize 见 \ref{sec:parityform}~节}};
  \node[box, below right=0.7cm and 0.2cm of drv] (ipopt) {\srcpath{solve\_parity\_with\_ipopt}\\[-2pt]
        {\footnotesize src/optimal\_power\_flow/ac\_opf.cpp:solve\_parity\_with\_ipopt（\srcpath{HACDCPF\_HAVE\_IPOPT} 门控）}};
  \node[box, below=2.0cm of drv] (core) {Parity IPM 核\\
        \srcpath{solve\_primal\_dual\_ipm}\\[-2pt]{\footnotesize src/optimal\_power\_flow/parity\_ipm.cpp:solve\_primal\_dual\_ipm}};
  \node[box, below left=0.7cm and 0.2cm of core] (kkt) {KKT 线性代数\\
        稠密 LU / UMFPACK / KLU\\/ Eigen SparseLU（MUMPS 被门控）\\[-2pt]
        {\footnotesize src/optimal\_power\_flow/parity\_ipm.cpp:backend\_order}};
  \node[box, below right=0.7cm and 0.2cm of core] (res) {\fld{IPMResult}\\
        收敛标志/残差/状态串\\[-2pt]{\footnotesize include/hacdcpf/optimal\_power\_flow/native\_ipm\_solver.hpp:IPMResult}};
  \draw[arr] (api) -- (drv);
  \draw[arr] (drv) -- (prob);
  \draw[arr] (drv) -- (ipopt);
  \draw[arr] (prob) -- (core);
  \draw[arr] (drv) -- node[lab, right=2pt]{NativeIPM/Auto} (core);
  \draw[arr] (core) -- (kkt);
  \draw[arr] (core) -- (res);
  \draw[arr, dashed] (ipopt) -- node[lab, above=1pt]{Auto 回退} (res);
\end{tikzpicture}}
\caption{Parity IPM 核的调用链与门控位置。虚线为仅当本核未收敛且
\fld{allow\_fallback} 时启用的 Ipopt 替代路径（\ref{subsec:pipm-ipopt}~节）。}
\label{fig:pipm-chain}
\end{figure}

\subsubsection*{选项与结果结构体}
\compmeta{IPMOptions}{\srcpath{include/hacdcpf/optimal\_power\_flow/native\_ipm\_solver.hpp:IPMOptions}}
\fld{IPMOptions} 承载迭代上限、三类容忍度、正则化基数、fraction-to-boundary
因子 $\tau$ 与可选的完整热启动状态（原始点、等式/不等式乘子、松弛）；
全字段见 \ref{subsec:pipm-params}~节参数表。如实说明两点：其一，
\fld{tol\_complementarity} 与 \fld{regularization} 两个字段在本核实现中
\textbf{从未被读取}（:solve\_primal\_dual\_ipm 内仅出现 \fld{tol\_primal}、\fld{tol\_dual}、
\fld{alpha\_max}、\fld{max\_iter} 与 \fld{verbose}），互补容忍度由
\fld{tol\_primal} 派生、正则化梯级为硬编码常数（见
\ref{subsec:pipm-conv}、\ref{subsec:pipm-kkt}~节）；其二，\fld{tol\_dual}
不进入收敛判据本身，仅用于滤波器线搜索的平稳性守卫（:solve\_primal\_dual\_ipm）。

\compmeta{IPMResult}{\srcpath{include/hacdcpf/optimal\_power\_flow/native\_ipm\_solver.hpp:IPMResult}}
\fld{IPMResult} 返回收敛标志、迭代数、三类归一化残差
（\fld{primal\_inf}/\fld{dual\_inf}/\fld{complementarity}）、初值残差、
热启动标记、线性代数计数（\fld{factorization\_calls}、
\fld{linear\_solve\_calls}、\fld{accepted\_steps}、\fld{rejected\_steps}、
\fld{symbolic\_analyze\_calls}、\fld{backend\_escalations}、
\fld{scaling\_rebuilds}）、英文状态串 \fld{status}、实际 KKT 后端标签
\fld{linear\_solver}（\fld{dense\_lu}/\fld{sparse\_umfpack}/
\fld{sparse\_klu}/\fld{sparse\_eigen\_lu}/\fld{sparse}，:solve\_primal\_dual\_ipm）
、Phase I 前后违反量/dual-fit/预算/终止原因、Phase II 接受标志，
与最终原始-对偶状态 $(x,\lambda,\mu,z)$。状态串的可取值全集见
\ref{subsec:pipm-errors}~节的状态串表。

\subsection{障碍问题形式与中心路径参数}\label{subsec:pipm-barrier}

\subsubsection*{障碍问题与松弛化}
核求解的问题是 Parity 建模层装配的非线性规划（头文件问题声明
\srcpath{native\_ipm\_solver.hpp:solve\_primal\_dual\_ipm}）：
\begin{align*}
\min_{x}\; & f(x)\\
\text{s.t.}\; & g(x) = 0,\quad h(x) \le 0,\quad x_{\min} \le x \le x_{\max}
\end{align*}
其中 $x\in\mathbb{R}^{n}$ 为 \srcpath{Problem::vidx} 定义的全空间变量
（$n=$~\fld{n\_total}），$g:\mathbb{R}^n\to\mathbb{R}^{m_{eq}}$ 为等式行
（$m_{eq}=$~\fld{n\_eq\_total}，含 AC/DC 功率平衡、换流器平衡、参考角与
直流参考电压），$h:\mathbb{R}^n\to\mathbb{R}^{m_{nl}}$ 为非线性不等式
（$m_{nl}=$~\fld{n\_ineq\_nonlin}，线路容量、VSC 视在功率/电流/调制窗口等）。
核在入口处把变量上下界改写为两行一组的一次不等式并入 $h$：对每个有限下界
列 $c$ 追加 $x_{\min,c}-x_c\le 0$（雅可比行仅含 $-1$），对每个有限上界列
追加 $x_c-x_{\max,c}\le 0$（仅含 $+1$）（\srcpath{assemble\_inequalities}，
界列扫描 :solve\_primal\_dual\_ipm）。合并后的不等式总数记
$m=n_{iq}=m_{nl}+n_{lb}+n_{ub}$；若 $n_{iq}=0$ 直接抛异常
（:solve\_primal\_dual\_ipm）。引入松弛 $z>0$ 与不等式乘子 $\mu>0$ 后，核迭代求解
扰动 KKT 系统
\begin{align*}
\nabla f(x) + J_g(x)^{\top}\lambda + J_h(x)^{\top}\mu &= 0\\
g(x) &= 0\\
h(x) + z &= 0\\
Z\mu &= \gamma\, e
\end{align*}
（\srcpath{native\_ipm\_solver.hpp:solve\_primal\_dual\_ipm} 文档注释与 :solve\_primal\_dual\_ipm 的
右端组装一致；$Z=\mathrm{diag}(z)$，$e$ 为全 1 向量，$\gamma>0$ 为当前
障碍目标）。建模层对目标梯度作工程缩放：当
$\lVert\nabla f\rVert_\infty>100$ 时取
$\omega_f=\min(1,\,100/\lVert\nabla f\rVert_\infty)$ 统一缩小目标及其
对角 Hessian，把拉氏梯度与最优乘子压回 $O(1)$ 量级（:solve\_primal\_dual\_ipm、:apply\_obj\_scale\_hessian（lambda），参考
Ipopt \fld{nlp\_scaling\_method} 的做法）；下文 $\tilde f=\omega_f f$，
$\tilde\nabla f$、$\nabla^2\tilde f$ 同义。

\subsubsection*{初始点与初值规则}
原始初值优先取调用方热启动点：逐变量投影到界内，并保留极小内点余量
（:solve\_primal\_dual\_ipm）$\varepsilon=\min(\,10^{-6}\max(1,\mathrm{width}),$
$0.49\,\mathrm{width}\,)$；无热启动时由建模层 \srcpath{build\_initial\_point}
给点
（\srcpath{formulation.hpp:build\_initial\_point}）。冷启动松弛取
$z_i=\max(|h_i(x_0)|,\,1)$（$h_i\ge0$）或 $\max(-h_i,\,10^{-2})$
（$h_i<0$），乘子取 $\mu_i=\max(1/z_i,\,10^{-2})$ 以维持
$\mu_i z_i\approx1$；等式乘子 $\lambda=0$（:solve\_primal\_dual\_ipm）。热启动四元组
仅在维度、有限性与正性检查全部通过时整体采用（:solve\_primal\_dual\_ipm）。
$\tau$ 取 \fld{alpha\_max}（:solve\_primal\_dual\_ipm）。

\subsubsection*{有界 Phase I 原--对偶起点构造}
在没有兼容完整四元组且 \fld{enable\_phase\_one=true} 时，核先运行局部约束
Newton 恢复。标准 AC 潮流分区把电压作为 state/basic 列；发电机母线以一个
$Q_g$ 替代对应 $V_m$，参考母线再以一个 $P_g$ 完成方阵。Hybrid AC/DC 的
$V_{dc}$、换流器/DC/DC/能量路由器平衡列按同一规则补齐。若当前没有违反的
非线性不等式且方阵可分解，求
\begin{equation}
 J_B(x_k)\Delta x_B=-g(x_k),\qquad \Delta x_F=0;
\end{equation}
否则把所有 $h_i(x_k)>\varepsilon_p$ 行追加到残差，对未固定在界上的列做一次
变量区间缩放后的 SparseQR。该构造对应 Nocedal--Wright（2006）第 7.5、
11.1 与 16.5 节的缩放约束 Newton/基变量思想；线搜索按 Deuflhard（2011）
第 2.2--2.3 节只接受原 pu 坐标最大违反量严格下降的步。

Hybrid AC/DC 在上述全模型恢复前先利用 DC 电导块做一次结构步。对换流器行
$g_c$ 与 DC 平衡行 $g_d$ 的线性化写成
\begin{align}
J_{dv}\Delta V_{dc}+J_{dp}\Delta P_{dc}&=-g_d,\\
J_{cp}\Delta P_{dc}+J_{ca}\Delta P_{ac}&=-g_c.
\end{align}
先由第二式消去局部 $\Delta P_{dc}$，得到只含
$[\Delta V_{dc},\Delta P_{ac}]$ 的 Schur 系统；DC reference 行直接追加。
这一步只分解约 $n_{dc}+n_{conv}$ 的稀疏块，不形成 OPF Hessian/KKT。接受准则
要求 DC balance/converter/reference 子块无穷范数严格下降，且全模型原坐标
违反量不增；全局 admission 仍然独立生效，因此局部 DC 改善不能强迫远离
可行域的 case300 进入昂贵全模型 QR。诊断字段
\fld{phase\_one\_structural\_*} 与 \fld{phase\_one\_structure} 区分该步、
后续 state/basic Newton 和普通回退。

Phase I 与 Phase II 的默认耦合参数为
$\mu_0=0.1$、$\eta_a=1$、$\eta_p=0.1$、$\eta_c=0.5$。只有初始违反量满足
$\|r_p(x_0)\|_\infty\le\eta_a\mu_0$ 才进入局部恢复；否则以
\fld{outside-phase-one-admission} 零分解退出，且不改变 Phase II 冷起点。
恢复在 $\epsilon_p=\max(\mathrm{tol}_p,\eta_p\mu_0)$ 停止，而不是追到最终
精度。令 $v=\max(0,\max_i h_i)$，松弛下限取
$s_{\min}=\max(2\times10^{-10},(\epsilon_p-v)/2)$，于是
$h_i+s_i\le(\epsilon_p+v)/2\le\epsilon_p$；再置
$\nu_i=\mu_0/s_i$，避免仅用机器下限造成巨大乘子，同时精确保持中心乘积。
移交还要求 $\max_i|s_i\nu_i/\mu_0-1|\le\eta_c$。

Phase I 只有迭代、稀疏分解、回溯与协作式墙钟预算，不含障碍序列、filter、
Hessian、二阶校正或弹性变量。原始点认证后只做一次 $J_g^T\lambda$ 对偶拟合，
\fld{phase\_one\_dual\_fit\_residual} 报告 basic 列 stationarity（或 SparseQR
正规残差）的缩放范数，保留自由控制列的 reduced gradient 给 Phase II。
完整 continuation state 直接绕过 Phase I。\fld{phase\_one\_primal\_feasible}
仍专指达到最终容差；\fld{phase\_one\_in\_handoff\_corridor} 才表示达到上述
有限精度移交邻域。若 primal/dual/centrality 证书任一失败，
Phase II 精确恢复进入 Phase I 之前的原始点以及普通冷启动 slack/dual；因此
\fld{phase\_two\_start\_accepted=false} 的试探轨迹不会改变最优性求解。
这一区分使 Phase I 保持 warm-start producer，而不是第二套 infeasible-start IPM。

\subsubsection*{Prepared OPF session 与完整状态延续}
公开 RAII 类型 \srcpath{PreparedACOPFSession} 为重复 Parity ACOPF 保留
formulation、组件/行映射、稀疏 pattern 的 symbolic ordering，以及上一轮完整
$(x,\lambda,\mu,z)$。兼容键不是维数检查，而是覆盖变量块、等式/不等式块、
有限 bound-row 顺序、reference row、受限设备集合、energy-router port 顺序与
$Y_{bus}/G_{dc}$ 非零 pattern 的 64 位签名；LCC 的稳定 ID、AC/DC 端点、解析后的
换流母线以及 energy-router port 稳定 ID 同样进入语义映射审计。参数变化先重新投影 canonical numeric
\srcpath{SolverData}；若结构审计通过，只刷新负荷/ZIP、固定注入、目标/界与矩阵
数值，保留 OPF maps 和 symbolic ordering。拓扑、服务状态、PV/PQ/reference、
constraint-family 或 row membership 改变时完整 rebuild，并在该次求解禁止 opaque
continuation state。linear backend 还对压缩 KKT 的全部行列坐标计算 pattern 指纹；
维数和非零元个数相同但坐标不同仍强制重新 symbolic analyze。

四块状态必须同时存在、签名相同、维数正确、数值有限且 slack/inequality dual
为正，随后仍通过本节 Phase-II residual/centrality 审计；任一条件失败时只允许
公共物理变量映射，不允许“维数碰巧相同”的 dual row 复用。当前 SuiteSparse
adapter 与 MIPSolvers MUMPS wrapper 只支持复用 symbolic 后重新 numeric factorize，
没有可按矩阵漂移、scaled backward residual 与 refinement 独立审计的跨求解
numeric-factor 对象。因此 \fld{prepared\_numeric\_refactor} 默认关闭；即使开启也
返回 \fld{unsupported} 并做新 numeric factorization，不虚报 numeric reuse。
AC-PF adjoint dual seed 未实现，因为 reduced PV/PQ Jacobian 与 full-space OPF
equality/bound multiplier 布局不兼容，且未保留 PF factor 时仍需额外分解。

\subsubsection*{预测--校正与中心路径参数更新}
每次迭代先对 $z,\mu$ 施加 $10^{-12}$ 下限保护并构成障碍比
$\Sigma=\mathrm{diag}(\mu_i/z_i)$（:solve\_primal\_dual\_ipm），随后执行
Mehrotra 预测--校正。预测步以 $\gamma=0$ 解牛顿系统得仿射方向
$(\Delta x^{\mathrm{aff}},\Delta\lambda^{\mathrm{aff}},
\Delta z^{\mathrm{aff}},\Delta\mu^{\mathrm{aff}})$，并以 $\tau=1$ 探测
仿射步长（凝集与增广形式均见 :solve\_primal\_dual\_ipm）：
\begin{align*}
\alpha_p^{\mathrm{aff}} &= \min\{\mathrm{ftb}(z,\Delta z^{\mathrm{aff}},1),\,1\},
&
\alpha_d^{\mathrm{aff}} &= \min\{\mathrm{ftb}(\mu,\Delta\mu^{\mathrm{aff}},1),\,1\}
\end{align*}
对偶间隙与仿射间隙为（凝集与增广均见 :solve\_primal\_dual\_ipm）
\begin{align*}
\bar\mu &= \frac{z^{\top}\mu}{n_{iq}},
&
\bar\mu_{\mathrm{aff}} &=
\frac{(z+\alpha_p^{\mathrm{aff}}\Delta z^{\mathrm{aff}})^{\top}
      (\mu+\alpha_d^{\mathrm{aff}}\Delta\mu^{\mathrm{aff}})}{n_{iq}}
\end{align*}
定心参数取立方律并截断（凝集与增广均见 :solve\_primal\_dual\_ipm）：
\begin{align*}
\sigma &= \left(\mathrm{clamp}\!\left(
  \frac{\bar\mu_{\mathrm{aff}}}{\bar\mu+10^{-16}},\,0,\,1\right)\right)^{3},
&
\gamma &= \mathrm{clamp}\!\left(\sigma\,\bar\mu,\;10^{-14},\,10^{4}\right)
\end{align*}
校正步右端补入 Mehrotra 二阶交叉项
$c_i=\mathrm{clamp}(\Delta\mu_i^{\mathrm{aff}}\Delta z_i^{\mathrm{aff}},
\,\pm 0.1\gamma)$（:solve\_primal\_dual\_ipm）。凝集形式的校正右端为
（:solve\_primal\_dual\_ipm）
\begin{align*}
N_{\mathrm{corr}} &= \nabla\tilde f + J_g^{\top}\lambda + J_h^{\top}\mu
+ J_h^{\top}\!\left((\mu\circ h + \gamma e - c)\oslash z\right)
\end{align*}
其中 $\circ$、$\oslash$ 为逐分量乘/除；校正方向的松弛与乘子恢复为
（:solve\_primal\_dual\_ipm）
\begin{align*}
\Delta z &= -(h+z) - J_h\,\Delta x,
&
\Delta\mu &= -\left(\mu + \Sigma\,\Delta z + (c-\gamma e)\oslash z\right)
\end{align*}
此后可追加至多 2 次 Gondzio 额外校正：对试步互补对
$\tilde z_i\tilde\mu_i$ 低于 $0.1\gamma$ 的分量以零等式残差右端再解一次
同分解方程并累加方向（凝集与增广均见 :solve\_primal\_dual\_ipm，阈值
$10^{-12}\gamma$）。增广形式另有两条防线：当 fraction-to-boundary 步长
被压到 $10^{-9}$ 以下时判定方向数值爆炸，把 $\gamma$ 提升到
$\max(\gamma,\,0.5\bar\mu)$ 并清零交叉项后重解校正，至多 2 次
（重新定心，:solve\_primal\_dual\_ipm）；滤波器全拒时还可以 $\gamma\leftarrow\max(\gamma,\bar\mu)$
的纯定心步整体重跑一次线搜索（中心性恢复重试，:solve\_primal\_dual\_ipm）。

环境变量 \srcpath{HACDCPF\_OPF\_MU\_STRATEGY=abo}（仅增广形式生效，
:solve\_primal\_dual\_ipm）把立方律换成 Armand--Benoist--Orban 动态 $\mu$：以
$\mu^{+}=\mu_{\mathrm{dyn}}^{2}/\mu_0$ 为更新函数、
\begin{align*}
\gamma &= \mathrm{clamp}\!\left(
\mu_{\mathrm{dyn}} + \alpha_{\mathrm{aff}}(\mu^{+}-\mu_{\mathrm{dyn}}),
\;10^{-14},\,10^{4}\right)
\end{align*}
其中 $\alpha_{\mathrm{aff}}=\min(\alpha_p^{\mathrm{aff}},
\alpha_d^{\mathrm{aff}},1)$，$\mu_0$ 为首个迭代的 $\bar\mu$
（:solve\_primal\_dual\_ipm）；实际采用的 $\gamma$（含重新定心修正）回写为下一迭代的
$\mu_{\mathrm{dyn}}$（:solve\_primal\_dual\_ipm）。默认 \fld{mehrotra}。

\subsection{KKT 系统组装与线性求解}\label{subsec:pipm-kkt}

\subsubsection*{两种牛顿系统形式}
核提供两种 KKT 组装形式（枚举 \fld{KktForm}；环境变量
\srcpath{HACDCPF\_OPF\_KKT\_FORM}，默认 \fld{condensed}，:kkt\_form\_preference）：
\begin{itemize}
  \item \textbf{凝集式（condensed，生产默认）}：把不等式乘子消入
        $(1,1)$ 块，
        \begin{align*}
        W &= \nabla^{2}_{xx}\tilde L + J_h^{\top}\Sigma J_h,
        &
        \begin{bmatrix} W+\delta I & J_g^{\top}\\ J_g & -\delta I\end{bmatrix}
        \begin{bmatrix}\Delta x\\ \Delta\lambda\end{bmatrix}
        &= -\begin{bmatrix} N\\ g\end{bmatrix}
        \end{align*}
        $W$ 的三重积按 $\mathrm{diag}(\Sigma)\cdot J_h$ 先行缩放再乘
        装配（:solve\_primal\_dual\_ipm），KKT 三元组装配见
        \srcpath{factor\_kkt\_sparse}（:factor\_kkt\_sparse）；$N$ 在预测步取
        $\tilde L_x + J_h^{\top}(\mu\circ h\oslash z)$（:solve\_primal\_dual\_ipm），
        校正步取 $N_{\mathrm{corr}}$。头文件注释指出凝集式条件数被
        平方（$\kappa\sim\sigma^2$，$\sigma=\mu/z$ 跨 $10^{\pm10}$），
        是 PEGASE 算例滤波器停滞的根因之一（:KktForm）。
  \item \textbf{增广式（augmented）}：保留不等式乘子为未知量，
        \begin{align*}
        \begin{bmatrix}
        \nabla^{2}_{xx}\tilde L+\delta_W I & J_g^{\top} & J_h^{\top}\\
        J_g & -\delta_C I & 0\\
        J_h & 0 & -ZM^{-1}
        \end{bmatrix}
        \begin{bmatrix}\Delta x\\ \Delta\lambda\\ \Delta\mu\end{bmatrix}
        = \begin{bmatrix}-\tilde L_x\\ -g\\ r_3\end{bmatrix}
        \end{align*}
        （块结构注释、装配与 $(3,3)$ 块 $-z_i/\mu_i$ 均见
        \srcpath{factor\_kkt\_sparse\_augmented}）。条件数对障碍比为
        线性，且对称不定分解可精确核对惯性：约化 Hessian 正定时惯性应为
        $(n,\;m_{eq}+n_{iq},\;0)$。
\end{itemize}
稠密路径（下文）与 $n_{iq}=0$ 时一律回落凝集式（:solve\_primal\_dual\_ipm）。

\subsubsection*{正则化与惯性控制}
凝集式采用渐进正则化：依次尝试
$\delta\in\{0,\;10^{-8}\lVert W\rVert,\;10^{-6}\lVert W\rVert,
\;10^{-4}\lVert W\rVert\}$（下限 $10^{-12}$），先分解再回代，任一失败
即升级，四级全败报 \texttt{"KKT factorization failed"}（:solve\_primal\_dual\_ipm）。
增广式实现 Wächter--Biegler $\delta_W$ 惯性校正：仅当 LDL$^{\top}$ 后端
（MUMPS）报告负主元数且不等于 $m_{eq}+n_{iq}$ 时，以
$\delta_W=10^{-8}\lVert\nabla^2_{xx}\tilde L\rVert$ 起步、每次 $\times8$、
上限 $10^{-2}\lVert\nabla^2_{xx}\tilde L\rVert$ 重分解，至多 8 次，末次
封顶后按 best-effort 接受（:solve\_primal\_dual\_ipm）；LU 类后端不报惯性
（\fld{negs}$=-1$）则\textbf{不加判别地接受未正则化分解}（:solve\_primal\_dual\_ipm）。
$\delta_W$ 热启动取 $\max(10^{-8}\lVert L_{xx}\rVert,\,0.5\,\delta_W^{prev})$
且只对 MUMPS 步启用，避免泄漏进 KLU 轨迹（:solve\_primal\_dual\_ipm）。$(2,2)$ 块正则
$\delta_C$ 默认 $10^{-8}\lVert L_{xx}\rVert$，可用
\srcpath{HACDCPF\_OPF\_DELTA\_C} 覆盖（:solve\_primal\_dual\_ipm）。如实说明：由于
MUMPS 在本检出处不可用（见下），$\delta_W$ 校正在实际构建中被静默绕过，
增广式的惯性保证随之失效。

\subsubsection*{平衡、分解与后端选择}
稀疏 KKT 缓存 \fld{SparseKKTCache}（:SparseKKTCache）按“模式不变”假设复用符号
分析：每次求解先对 KKT 做对称 Ruiz 平衡——迭代至多 5 轮、每轮把
$D_i$ 除以当前行 $\infty$-范数的平方根、行失衡偏差 $|\log r_{\max}|<10^{-3}$
即停（:compute\_ruiz\_scaling）——得到 $D\!KD$，此后求解按
$(D\!KD)\,y=D\,b$、$x=D\,y$ 进行（:kkt\_solve\_scaled）。平衡矩阵默认整个
求解过程只算一次并冻结（\srcpath{HACDCPF\_OPF\_RESCALE\_PER\_ITER=1}
恢复逐迭代重算，:factor\_assembled\_kkt）；注释记录了冻结缩放把 IEEE24 三区域扩展
算例从 640 迭代停滞修复为 36 迭代收敛的教训（:factor\_assembled\_kkt）。
回代自带至多 2 步迭代精化，且仅当校正\textbf{确实减小}残差时才采纳
（防数值奇异 KKT 上精化发散，:kkt\_solve\_scaled）。

线性后端按结构感知顺序选取（\srcpath{backend\_order}）：
\begin{itemize}
  \item 纯 AC 问题（无 DC 母线/支路、VSC、DC/DC、能量路由器，判定
        :solve\_primal\_dual\_ipm）：KLU $\to$ MUMPS $\to$ UMFPACK $\to$ Eigen
        SparseLU（:backend\_order）。KLU 的块三角分解匹配电网 KKT 的近块三角
        结构，注释给出 case2869 分解 0.006~s 对 MUMPS 0.015~s、端到端
        5.97~s 对 12.05~s 且目标值一致的实测（:backend\_order）；
  \item 混合 AC/DC 问题：MUMPS $\to$ UMFPACK $\to$ KLU $\to$ Eigen
        SparseLU（:backend\_order），理由是 $\delta_W$ 校正需要 LDL$^{\top}$
        惯性（:backend\_order）。
\end{itemize}
活动后端具有粘性；仅当 Eigen SparseLU 回代残差超过
$10^{-6}\lVert b\rVert_\infty$ 时置 \fld{solve\_degraded}，下次分解沿候选
序列升级一档（:kkt\_solve\_scaled、:factor\_assembled\_kkt）；UMFPACK/KLU 的告警残差不触发降级
（注释 :kkt\_solve\_scaled）。每个后端按 $(\dim,\mathrm{nnz})$ 缓存各自的符号分析，
模式不变即复用（:SparseKKTCache、:factor\_assembled\_kkt）。$n+m_{eq}\le512$ 的小系统改走
稠密 \fld{Eigen::PartialPivLU}（阈值 \fld{kDenseAutoKktDim}，:solve\_primal\_dual\_ipm；
稠密装配/分解 :factor\_kkt\_dense，同样带 2 步精化 :kkt\_solve\_dense）；$0\times0$
KKT 按 $\mathbb R^0$ 上唯一空解直接闭合，不调用 Eigen 分解。环境变量
\srcpath{HACDCPF\_OPF\_LINEAR\_SOLVER}（旧拼写
\srcpath{HACDCPF\_OPF\_SPARSE\_SOLVER} 兼容）可强制
\fld{dense}/\fld{mumps}/\fld{umfpack}/\fld{klu}/\fld{eigen}/\fld{auto}
（:linear\_solver\_preference\_env）。

\subsubsection*{稀疏库的编译期门控与回退}
UMFPACK/KLU 支持由 \srcpath{HACDCPF\_OPF\_HAVE\_UMFPACK} 与
\srcpath{HACDCPF\_OPF\_HAVE\_KLU} 门控（:SparseKKTCache、:backend\_order）：CMake 在
\srcpath{HACDCPF\_USE\_SUITESPARSE=ON}（默认 ON，\srcpath{CMakeLists.txt:HACDCPF\_USE\_SUITESPARSE}）
下优先链接 MIPSolvers 内置的 \fld{umfpack\_vendored}/\fld{klu\_vendored}
目标并定义两宏（\srcpath{CMakeLists.txt:HACDCPF\_OPF\_HAVE\_UMFPACK、:HACDCPF\_OPF\_HAVE\_KLU}），否则回退系统 SuiteSparse 探测
（\srcpath{CMakeLists.txt:HACDCPF\_SUITESPARSE\_INCLUDE\_DIR}）；两级都找不到时两个宏不定义，稀疏路径只剩
Eigen \fld{SparseLU}（COLAMD 排序，:SparseKKTCache、:sparse\_factorize\_active）。MUMPS 分支由
\srcpath{HACDCPF\_HAVE\_MUMPS} 门控（:SparseKKTCache、:backend\_order），文件头注释称该宏
“继承自 mipsolvers 目标的公共编译定义”；但本检出处 MIPSolvers 的 CMake
从不定义该宏（全仓检索仅在 parity\_ipm.cpp 内部出现），且其所引用的
\srcpath{mipsolvers::engine::MumpsSolver} 类在当前 MIPSolvers 头文件树中
不存在。因此\textbf{本检出处 MUMPS 恒被编译排除}：混合 AC/DC 问题的实际
后端序列为 UMFPACK $\to$ KLU $\to$ Eigen SparseLU，上文的 $\delta_W$
惯性校正与 case 注释中的 MUMPS 实测数字对当前构建均不适用。

\subsection{步长规则与全局化}\label{subsec:pipm-step}

\subsubsection*{fraction-to-boundary 与空步判定}
原始/对偶步长分别由 fraction-to-boundary 规则给出
（\srcpath{ftb}）：
\begin{align*}
\alpha(v,\Delta v,\tau) &= \min\!\left\{1,\;
\min_{i:\,\Delta v_i<0}\frac{-\tau\, v_i}{\Delta v_i}\right\}
\end{align*}
最终取 $\alpha_p=\max(\min(\mathrm{ftb}(z,\Delta z,\tau),1),10^{-10})$、
$\alpha_d$ 同理（:solve\_primal\_dual\_ipm）。当 $\min(\alpha_p,\alpha_d)<10^{-7}$
时判定该方向为空步（vacuous step）：滤波器的进步余量
$\eta\alpha$、$\gamma_\theta\theta$ 已数值失效，直接放弃线搜索进入
异常路径（:solve\_primal\_dual\_ipm；注释记录了 case13659 在 $\alpha\sim10^{-9}$
“被接受”步上随机游走数百迭代的事故）。

\subsubsection*{回溯线搜索与 $(\theta,\varphi)$ 滤波器}
线搜索对试步 $(x+\alpha_p\Delta x,\;z+\alpha_p\Delta z,\;
\mu+\alpha_d\Delta\mu,\;\lambda+\alpha_d\Delta\lambda)$ 重估全部非线性
残差并做接受测试，拒绝则 $\alpha\leftarrow0.5\alpha$，至多 14 次回溯
（\fld{kMaxBacktracks}/\fld{kBacktrack}，:solve\_primal\_dual\_ipm；主循环
\srcpath{run\_backtracks}（lambda））。接受准则有两层：
\begin{enumerate}
  \item \textbf{遗留三轴准则}：可行轴要求
        $\theta_t\le(1-\eta\bar\alpha_p)\theta$ 且对偶两轴不放大超 100 倍；
        平稳轴与互补轴分别在可行/平稳守卫
        （$\theta_t\le\max(10\varepsilon_p,1.02\theta)$、
        $\|L_x\|_t\le\max(10\varepsilon_d,1.02\|L_x\|)$）下要求对应量
        按 $(1-\eta\bar\alpha_d)$ 下降，$\eta=10^{-4}$（:eval\_and\_test\_trial（lambda））。
        注释指出互补轴守卫宽松（每步 2\%，复利）曾使 case2869pegase
        目标与平稳性虚胀约 3 倍，可用 \srcpath{HACDCPF\_OPF\_NO\_LEGACY=1}
        关闭（:solve\_primal\_dual\_ipm）；
  \item \textbf{$(\theta,\varphi)$ 滤波器}（Wächter--Biegler §3.2，Ipopt
        默认常数）：$\varphi_\gamma(x,z)=\tilde f(x)-\gamma\sum_i\ln z_i$，
        试点须相对当前对与\textbf{全部历史对}在
        $\theta$ 上占优（$\theta_t\le(1-\gamma_\theta)\theta$）或在
        $\varphi$ 上至少降 $\gamma_\varphi\theta$（支配记忆见
        :eval\_and\_test\_trial（lambda），历史入栈 :run\_backtracks（lambda））；近可行且方向为真实 $\varphi$ 下降
        时追加 Armijo 条件
        $\varphi_t\le\varphi_k+\eta_\varphi\bar\alpha\, d\varphi$
        （switching 条件与 Armijo 均见 :eval\_and\_test\_trial（lambda））。常数
        $\gamma_\theta=\gamma_\varphi=10^{-5}$、$\delta_s=1.0$、
        $s_\theta=1.1$、$\eta_\varphi=10^{-8}$、$\theta_{\min}=10^{-4}$
        （:solve\_primal\_dual\_ipm）；另有绝对天花板
        $\theta_{\max}=10^{4}\max(1,\theta_0)$（:solve\_primal\_dual\_ipm）。
\end{enumerate}
滤波器机制本身即非单调策略：接受只要求相对历史包络的进步，允许单轴
短期恶化，因此目标与残差沿轨迹均不单调；核另以 best-iterate 档案兜底
（:solve\_primal\_dual\_ipm）。如实说明：只有当 $n+m_{eq}\ge512$ 或
\srcpath{HACDCPF\_OPF\_FILTER\_ALL=1} 时线搜索才真正执行（:solve\_primal\_dual\_ipm）；
小 KKT 系统在 \srcpath{eval\_and\_test\_trial} 内无条件接受首个有限试步
（:eval\_and\_test\_trial（lambda）），即小算例走的是无阻尼全步牛顿轨迹，发散风险由后续异常路径
承担。诊断开关 \srcpath{HACDCPF\_OPF\_STRICT\_THETA=1} 把接受准则收紧为
纯可行下降（:solve\_primal\_dual\_ipm、:eval\_and\_test\_trial（lambda））。

\subsubsection*{二阶校正、中心性恢复与恢复阶段}
三个逐级升级的全局化补救（均仅增广式；凝集式失败后直接进入异常路径）：
\begin{enumerate}
  \item \textbf{二阶校正（SOC）}：首轮试步因可行性恶化被拒时，用同一
        分解把等式块右端换成 $-(\alpha_p g+g(x+\alpha_p\Delta x))$ 再解
        一次并尝试复合步，抵消 Maratos 效应的一阶残差上升，每迭代至多
        一次（:solve\_primal\_dual\_ipm 注释，实现 :run\_backtracks（lambda））；
  \item \textbf{中心性恢复重试}：滤波器全拒但首轮可行性尚可时，以
        $\gamma\leftarrow\max(\gamma,\bar\mu)$、交叉项清零重解校正方向
        并再跑一次线搜索（Gondzio--Grothey 式定心步，:solve\_primal\_dual\_ipm）；
  \item \textbf{恢复阶段}：以上均失败时，把 $(1,1)$ 块换成
        $\rho I$（$\rho=10^{-8}\lVert L_{xx}\rVert$）做一步 Gauss-Newton
        型可行性步，$\alpha$ 折半至多 8 次，接受首个
        $\theta$ 下降且对偶轴有界（$\le100$ 倍）的试步；下一主迭代用真实
        $L_{xx}$ 重新分解，等效于 Ipopt 恢复模式的有界版本
        （:solve\_primal\_dual\_ipm）。
\end{enumerate}

\subsection{收敛判据}\label{subsec:pipm-conv}

\subsubsection*{归一化残差度量}
核以四个归一化量度量收敛（\srcpath{compute\_convergence}）：
\begin{align*}
\theta &= \frac{\max(\lVert g\rVert_\infty,\,\max_i h_i^{+})}
               {1+\max(\lVert x\rVert_\infty,\lVert z\rVert_\infty)},
&
\theta_{\mathrm{raw}} &= \max(\lVert g\rVert_\infty,\,\max_i h_i^{+})\\
\kappa_d &= \frac{\lVert \nabla\tilde f+J_g^{\top}\lambda
            +J_h^{\top}\mu\rVert_\infty}
            {1+\max(\lVert\mu\rVert_\infty,\lVert\lambda\rVert_\infty)},
&
\kappa_c &= \frac{z^{\top}\mu/n_{iq}}{1+\lVert x\rVert_\infty}
\end{align*}
另有成本条件 $\kappa_f=|\tilde f-\tilde f_0|/(1+|\tilde f_0|)$（:compute\_convergence（lambda）），
但如实说明：$\kappa_f$ 仅用于日志输出，\textbf{不进入任何收敛判据}。

\subsubsection*{判据全表}
下表“典型范围”列填判据的不等式形式，“默认值”列填默认容忍度，“说明”列
末注源码行号（均为 \srcpath{parity\_ipm.cpp}）。
\begin{paramtable}
原始可行（归一化） & $\theta$ & — & $<\varepsilon_p=\fld{tol\_primal}$ & $10^{-6}$ & 驱动层取 $\max(\fld{feasibility\_tol},10^{-10})$（src/optimal\_power\_flow/ac\_opf.cpp:solve\_with\_parity\_ipm；:compute\_convergence（lambda）、:solve\_primal\_dual\_ipm）\\
原始可行（原始量纲） & $\theta_{\mathrm{raw}}$ & — & $<\varepsilon_p$ & $10^{-6}$ & 与归一化值同时要求，防大 $\lVert x\rVert$ 稀释（:compute\_convergence（lambda）、:solve\_primal\_dual\_ipm）\\
对偶可行（平稳性） & $\kappa_d$ & — & $<\varepsilon_p$ & $10^{-6}$ & 注意：比较的是 $\varepsilon_p$ 而非 $\varepsilon_d$；$\fld{tol\_dual}$ 只进滤波器守卫（:compute\_convergence（lambda）、:solve\_primal\_dual\_ipm、:eval\_and\_test\_trial（lambda））\\
互补间隙 & $\kappa_c$ & — & $<\max(10\varepsilon_p,\,2\times10^{-3})$ & $2\times10^{-3}$ & 大算例障碍病态下平均 $z\mu/n_{iq}$ 常停在 $10^{-3}$ 级，故设工程下限；\fld{tol\_complementarity} 字段未被读取（:solve\_primal\_dual\_ipm）\\
迭代上限 & — & — & $\fld{max\_iter}$ & 640 & 驱动层折算为内、外上限乘积 $80\times8$（src/optimal\_power\_flow/ac\_opf.cpp:solve\_with\_parity\_ipm；:solve\_primal\_dual\_ipm）\\
可接受级收敛 & — & — & $\theta,\theta_{\mathrm{raw}},\kappa_d<100\,\varepsilon_p$ 且 $\kappa_c$ 达标 & $100\times$ & 取当前点或 best-iterate 中更优者，状态串如实标注 acceptable（Ipopt \fld{Solved\_To\_Acceptable\_Level} 惯例）（:solve\_primal\_dual\_ipm）\\
目标停滞早停 & — & — & $\theta_{\mathrm{raw}}<10^{-4}$、$\kappa_c$ 达标、$|\Delta \tilde f|\le10^{-6}(1+|\tilde f|)$ 连续 20 次 & 20 次 & 早停后状态为 \texttt{objective-stagnant (dual-polish certification pending)}，交由对偶抛光认证（:solve\_primal\_dual\_ipm）\\
对偶最小二乘认证 & — & — & 终点 $\theta,\theta_{\mathrm{raw}},\kappa_c$ 达 acceptable 且抛光后 $\kappa_d<100\,\varepsilon_p$ & $100\times$ & 解 $\min_{\lambda,\mu}\tfrac12\lVert\tilde\nabla f+J_g^{\top}\lambda+J_h^{\top}\mu\rVert^2$ 的 KKT 系统（$(1,1)$ 块取 $I$），$\mu$ 截断非负；状态串 \texttt{converged (dual least-squares certified)}（:solve\_primal\_dual\_ipm）\\
\end{paramtable}

严格判据为 $\theta<\varepsilon_p$、$\theta_{\mathrm{raw}}<\varepsilon_p$、
$\kappa_d<\varepsilon_p$、$\kappa_c<\kappa_c^{tol}$ 四者同时成立
（:solve\_primal\_dual\_ipm 初判与迭代内判定，置 \fld{converged} 与状态
\fld{"converged"}）。对偶抛光只在 $\kappa_d$ 不劣于轨迹自身乘子时采纳
抛光乘子（:solve\_primal\_dual\_ipm）。

\subsection{异常路径与结果标记}\label{subsec:pipm-errors}

\subsubsection*{核内状态串全集}
\fld{IPMResult::status} 的全部可取值及其触发位置如下（默认初值
\texttt{"maximum iterations reached"}，:solve\_primal\_dual\_ipm；“说明”列末注源码行号）：
\begin{paramtable}
\fld{converged} & — & — & — & — & 严格判据满足（:solve\_primal\_dual\_ipm）\\
\texttt{converged (acceptable tolerance, best iterate)} & — & — & — & — & best-iterate 达 $100\times$ 容忍度（:solve\_primal\_dual\_ipm）\\
\texttt{converged (acceptable tolerance, final iterate)} & — & — & — & — & 终点达 $100\times$ 容忍度（:solve\_primal\_dual\_ipm）\\
\texttt{converged (dual least-squares certified)} & — & — & — & — & 对偶抛光认证通过（:solve\_primal\_dual\_ipm）\\
\texttt{objective-stagnant (dual-polish certification pending)} & — & — & — & — & 目标停滞窗口触发；随后若抛光认证失败则保持未收敛（:solve\_primal\_dual\_ipm）\\
\texttt{maximum iterations reached} & — & — & — & — & 迭代耗尽（:solve\_primal\_dual\_ipm 默认值）\\
\texttt{KKT factorization failed} & — & — & — & — & 正则化梯级/惯性校正全败（:solve\_primal\_dual\_ipm 凝集与增广两支）\\
\texttt{predictor back-solve failed} & — & — & — & — & 预测步回代失败（:solve\_primal\_dual\_ipm，仅增广）\\
\texttt{NaN in corrector step} & — & — & — & — & 校正方向出现非有限值（:solve\_primal\_dual\_ipm）\\
\texttt{NaN in corrector step} & — & — & — & — & 校正方向出现非有限值（:build\_kkt\_workspace、:solve\_economic\_dispatch）\\
\texttt{filter line search failed} & — & — & — & — & 回溯耗尽且恢复阶段失败（:solve\_primal\_dual\_ipm）\\
\end{paramtable}
两条硬异常在入口处抛出 \fld{std::runtime\_error}：问题规模非法
（$n\le0$ 或 $m_{eq}<0$，:solve\_primal\_dual\_ipm）与装配后无不等式（:solve\_primal\_dual\_ipm）。

\subsubsection*{向 ACOPFResult 的映射（诚实口径）}
驱动层把核结果翻译为公共结果（src/optimal\_power\_flow/ac\_opf.cpp:solve\_with\_parity\_ipm）：
\begin{itemize}
  \item \fld{converged} 与 \fld{status}：收敛时
        \texttt{"converged (<backend>)"}（backend 如
        \fld{parity\_ipm:sparse\_klu}、\fld{ipopt\_filter\_linesearch}），
        Ipopt 可接受残差另加 \texttt{[acceptable residuals]}；未收敛时
        \texttt{"not converged: <IPMResult::status>"}，核的 acceptable/抛光
        语义透过状态串保留（src/optimal\_power\_flow/ac\_opf.cpp:solve\_with\_parity\_ipm）；
  \item 计数与残差进 \fld{ACOPFProfiling}：线性后端标签、分解/回代/
        接受/拒绝步数、符号分析次数、后端升级次数、缩放重建次数、热启动
        标记与初值残差（src/optimal\_power\_flow/ac\_opf.cpp:solve\_with\_parity\_ipm；结构体
        \srcpath{opf\_result.hpp:ACOPFProfiling}）；
  \item 原始解不存在或含非有限值时 \fld{objective} 与
        \fld{max\_constraint\_violation} 置 $+\infty$，跳过全部状态提取
        （src/optimal\_power\_flow/ac\_opf.cpp:solve\_with\_parity\_ipm）；
  \item 求解后在终态重估各族残差写入 \fld{profiling} 的
        \fld{max\_ac\_p\_balance\_residual\_pu} 等审计字段
        （src/optimal\_power\_flow/ac\_opf.cpp:solve\_with\_parity\_ipm）；未收敛时由 \srcpath{solve\_ac\_opf}
        汇总 \fld{infeasibility\_hints}（状态串、孤岛无松弛等，
        src/optimal\_power\_flow/ac\_opf.cpp:solve\_ac\_opf\_impl）；
  \item \fld{lmp\_valid} 仅在等式乘子长度覆盖 AC 平衡行时置真并注明
        来源，否则写明不可用原因（Ipopt 适配器不返回约束乘子时 LMP
        不可得，src/optimal\_power\_flow/ac\_opf.cpp:solve\_with\_parity\_ipm）；
  \item 模型覆盖声明走 \fld{converter\_model\_scope}/\fld{ValidityFlags}
        与 \fld{model\_limitations}（路径标签按实际启用的约束族拼接，
        LCC 准稳态、触发角/熄弧角求解后审计越限等逐条列出，
        src/optimal\_power\_flow/ac\_opf.cpp:solve\_with\_parity\_ipm），口径与 \ref{sec:overview}~节
        的诚实结果约定一致。
\end{itemize}

\subsection{嵌入式 Ipopt 替代路径}\label{subsec:pipm-ipopt}

\subsubsection*{编译期与运行期门控}
Parity 问题同时可交由嵌入式 Ipopt 求解（适配器
\srcpath{solve\_parity\_with\_ipopt}，src/optimal\_power\_flow/ac\_opf.cpp:solve\_parity\_with\_ipopt）：Parity
\fld{Problem} 暴露的目标、梯度、等式/不等式残差与雅可比、拉氏 Hessian
回调与 \fld{engine::NLPModel} 一一对应，界约束直接作为变量界传入
（src/optimal\_power\_flow/ac\_opf.cpp:build\_parity\_nlp\_input 注释）。门控分两层：
\begin{itemize}
  \item 编译期 \srcpath{HACDCPF\_HAVE\_IPOPT}（src/optimal\_power\_flow/ac\_opf.cpp:solve\_parity\_with\_ipopt）：
        平台选项 \srcpath{HACDCPF\_ENABLE\_IPOPT} 在 Apple 与 Windows 默认 ON、
        其余平台默认 OFF（\srcpath{CMakeLists.txt:HACDCPF\_ENABLE\_IPOPT}，任何平台均可显式开启）；
        Windows 默认选择 MIPSolvers 的 \fld{pardisomkl} 后端并要求本地静态
        sequential oneMKL 包完整，
        并强制透传为 \fld{MIPSOLVERS\_BUILD\_LOCAL\_IPOPT}；MIPSolvers 侧仅在以内嵌
        Ipopt 源码构建出 \fld{ipopt\_local} 时置 \fld{MIPSOLVERS\_HAVE\_IPOPT} 并透传
        \srcpath{HACDCPF\_HAVE\_IPOPT=1}（\srcpath{MIPSolvers/cmake/Dependencies.cmake:MIPSOLVERS\_HAVE\_IPOPT}、
        \srcpath{MIPSolvers/CMakeLists.txt:HACDCPF\_HAVE\_IPOPT}）——开启但内嵌
        \texttt{ipopt/} 源码缺失时在 MIPSolvers 侧 \fld{FATAL\_ERROR}（必须用仓库内
        Ipopt 源码，不用系统 Ipopt）。未编译入时适配器返回
        \fld{available=false}、状态 \texttt{"Ipopt not compiled in"}
        （src/optimal\_power\_flow/ac\_opf.cpp:solve\_parity\_with\_ipopt）；
  \item 运行期 \srcpath{HACDCPF\_DISABLE\_IPOPT\_OPF}：设置即强制不可用，
        用于后端隔离与诊断（src/optimal\_power\_flow/ac\_opf.cpp:solve\_parity\_with\_ipopt）。
\end{itemize}

\subsubsection*{与核的衔接方式}
内层求解器为 \fld{Auto} 且关闭目标同伦时，若启用了混合潮流热启动
（\fld{ac\_pf\_warm\_start}），先用 Ipopt 解、其未收敛但给出有限原始点时
再以完整四元组热启动本核精化（src/optimal\_power\_flow/ac\_opf.cpp:solve\_with\_parity\_ipm）；默认 \fld{Auto}
分支为本核先行、未收敛且 \fld{allow\_fallback} 时回退 Ipopt
（src/optimal\_power\_flow/ac\_opf.cpp:solve\_with\_parity\_ipm）；显式 \fld{Ipopt} 后端在 Ipopt 不可用时回落
本核并在后端标签中注明 \texttt{(ipopt unavailable)}（src/optimal\_power\_flow/ac\_opf.cpp:solve\_with\_parity\_ipm）。
如实说明两条口径差异：其一，Ipopt 报
\fld{Maximum\_Iterations\_Exceeded} 但原始可行严格达标、对偶与互补落在
$1000\times$ 容忍带（互补另取 $10^{-2}$ 工程下限）时仍按收敛接受，状态
串保留原终止码（src/optimal\_power\_flow/ac\_opf.cpp:solve\_parity\_with\_ipopt）；其二，适配器不返回
界乘子，驱动层以 $2\times10^{-12}$ 内点下限重构 $\mu,z$ 以满足核状态
格式（src/optimal\_power\_flow/ac\_opf.cpp:solve\_parity\_with\_ipopt），因此 Ipopt 路径的 LMP 不可用
（src/optimal\_power\_flow/ac\_opf.cpp:solve\_with\_parity\_ipm）。

\subsection{调谐参数表}\label{subsec:pipm-params}

\subsubsection*{IPMOptions 字段}
\begin{paramtable}
\fld{verbose} & — & — & true/false & false & 逐迭代调试日志（:IPMOptions，:solve\_primal\_dual\_ipm）\\
\fld{tol\_primal} & $\varepsilon_p$ & — & $10^{-8}$--$10^{-5}$（经验值） & $10^{-6}$ & 原始可行与平稳性判据容忍度；驱动层取 $\max(\fld{feasibility\_tol},10^{-10})$（\srcpath{native\_ipm\_solver.hpp}:IPMOptions，src/optimal\_power\_flow/ac\_opf.cpp:solve\_with\_parity\_ipm）\\
\fld{tol\_dual} & $\varepsilon_d$ & — & $10^{-8}$--$10^{-5}$（经验值） & $10^{-6}$ & 不进收敛判据；仅滤波器平稳守卫 $\max(10\varepsilon_d,1.02\kappa_d)$（\srcpath{native\_ipm\_solver.hpp}:IPMOptions，src/optimal\_power\_flow/parity\_ipm.cpp:solve\_primal\_dual\_ipm）\\
\fld{tol\_complementarity} & $\varepsilon_c$ & — & — & $10^{-6}$ & \textbf{本核未读取}；互补阈值固定为 $\max(10\varepsilon_p,2\times10^{-3})$（\srcpath{native\_ipm\_solver.hpp}:IPMOptions，src/optimal\_power\_flow/parity\_ipm.cpp:solve\_primal\_dual\_ipm）\\
\fld{regularization} & $\delta_0$ & — & — & $10^{-8}$ & \textbf{本核未读取}；正则化梯级硬编码（\srcpath{native\_ipm\_solver.hpp}:IPMOptions，src/optimal\_power\_flow/parity\_ipm.cpp:solve\_primal\_dual\_ipm）\\
\fld{alpha\_max} & $\tau$ & — & 0.9--0.9999（经验值） & 0.95 & fraction-to-boundary 因子；驱动层把 \fld{interior\_fraction} 经 clamp 夹到 $[0.5,\,0.9999]$，默认 0.995（\srcpath{native\_ipm\_solver.hpp}:IPMOptions，src/optimal\_power\_flow/ac\_opf.cpp:solve\_with\_parity\_ipm）\\
\fld{verbose} & — & — & true/false & false & 逐迭代调试日志（\srcpath{native\_ipm\_solver.hpp}:IPMOptions，src/optimal\_power\_flow/parity\_ipm.cpp:solve\_primal\_dual\_ipm）\\
\fld{enable\_phase\_one} & — & — & true/false & true & 无兼容完整状态时运行有界 primal/dual 起点构造\\
\fld{phase\_one\_time\_limit\_ms} & — & ms & $>0$ 协作截止；$\le0$ 不限 & 5000 & 在稀疏分解之间检查，单次分解不可中断\\
\fld{phase\_one\_max\_iterations} & — & — & $\ge0$ & 12 & 约束 Newton 迭代硬上限\\
\fld{phase\_one\_max\_factorizations} & — & — & $\ge0$ & 14 & primal 恢复与一次 dual fit 共用的数值分解硬上限\\
\fld{phase\_one\_max\_backtracks} & — & — & $\ge0$ & 12 & 每次 Newton 步的回溯上限\\
\fld{phase\_one\_dispatch\_dual\_min\_improvement} & — & — & 0--1 & 0.05 & raw 与 normalized dual residual 必须同时达到的改善比例（:IPMOptions，:solve\_primal\_dual\_ipm）\\
\fld{phase\_one\_admission\_mu\_factor} & — & — & $\ge0$ & 1.0 & Phase I 准入要求 violation $\le$ 因子$\times\mu_I$，否则记 \texttt{outside-phase-one-admission}；非有限回退 1.0（:inf\_norm，:elapsed\_ms（lambda））\\
\fld{phase\_one\_primal\_mu\_factor} & — & — & $\ge0$ & 0.1 & 交接 primal 容差 $\max(\varepsilon_p,\text{因子}\times\mu_I)$；非有限回退 0.1（:inf\_norm，:restore\_phase\_one\_primal）\\
\fld{phase\_one\_centrality\_tolerance} & $\eta_c$ & — & $\ge0$ & 0.5 & 中心性门槛（$\max_i|z_i\mu_i/\mu_0-1|$ 判据）；非有限回退 0.5（:ftb，:initialize\_phase\_one\_duals）\\
\fld{phase\_one\_dispatch\_dual\_predictor} & — & — & true/false & false & 实验性零分解连通分量 dispatch-price dual 候选（:ftb，:apply\_obj\_scale\_hessian（lambda））\\
\fld{phase\_one\_dispatch\_dual\_min\_improvement} & — & — & 0--1 & 0.05 & raw 与 normalized dual residual 必须同时达到的改善比例（:ftb，:apply\_obj\_scale\_hessian（lambda））\\
\fld{prepared\_state} & — & — & 指针/null & null & 跨求解保有 KKT 符号分析的所有者；\fld{layout\_signature} 变化时调用方须 \texttt{reset}（:IPMOptions，:solve\_primal\_dual\_ipm）\\
\fld{prepared\_numeric\_refactor} & — & — & true/false & false & 实验性数值重分解请求；当前后端无独立可审计 refactor 对象，fail-closed 报 \texttt{unsupported} 并总是新算数值因子（:IPMOptions，:solve\_primal\_dual\_ipm）\\
\fld{prepared\_numeric\_max\_relative\_drift} & — & — & $\ge0$ & $10^{-3}$ & 数值重分解允许的矩阵值相对漂移上限（:IPMOptions；当前后端不消费，见上行）\\
\fld{prepared\_numeric\_backward\_error\_tolerance} & — & — & $\ge0$ & $10^{-8}$ & 数值重分解的后向误差审计门槛（:IPMOptions；当前后端不消费，见上行）\\
\fld{slack\_start} & $z_0$ & — & — & nullptr & 热启动松弛，须严格为正（:IPMOptions，:solve\_primal\_dual\_ipm）\\
\fld{equality\_dual\_start} & $\lambda_0$ & — & — & nullptr & 热启动等式乘子，与下两项同时通过检查才启用（\srcpath{native\_ipm\_solver.hpp}:IPMOptions，src/optimal\_power\_flow/parity\_ipm.cpp:solve\_primal\_dual\_ipm）\\
\fld{inequality\_dual\_start} & $\mu_0$ & — & — & nullptr & 热启动不等式乘子，须严格为正（\srcpath{native\_ipm\_solver.hpp}:IPMOptions，src/optimal\_power\_flow/parity\_ipm.cpp:solve\_primal\_dual\_ipm）\\
\fld{slack\_start} & $z_0$ & — & — & nullptr & 热启动松弛，须严格为正（\srcpath{native\_ipm\_solver.hpp}:IPMOptions，src/optimal\_power\_flow/parity\_ipm.cpp:solve\_primal\_dual\_ipm）\\
\end{paramtable}

\subsubsection*{环境变量开关}
\begin{paramtable}
\srcpath{HACDCPF\_OPF\_LINEAR\_SOLVER} & — & — & dense/mumps/\allowbreak umfpack/klu/\allowbreak eigen/auto & auto & 强制 KKT 线性后端（旧拼写 \srcpath{HACDCPF\_OPF\_SPARSE\_SOLVER} 兼容）（:linear\_solver\_preference\_env）\\
\srcpath{HACDCPF\_OPF\_KKT\_FORM} & — & — & condensed/\allowbreak augmented & condensed & 牛顿系统形式；condensed 为生产默认，augmented 供刚性模型（:kkt\_form\_preference）\\
\srcpath{HACDCPF\_OPF\_MU\_STRATEGY} & — & — & mehrotra/abo & mehrotra & 中心路径参数策略（ABO 仅增广式）（:solve\_primal\_dual\_ipm）\\
\srcpath{HACDCPF\_OPF\_DELTA\_C} & $\delta_C$ & — & $10^{-10}$--$10^{-6}$（经验值） & $10^{-8}\lVert L_{xx}\rVert$ & 增广式 $(2,2)$ 块正则（绝对值）（:solve\_primal\_dual\_ipm）\\
\srcpath{HACDCPF\_OPF\_INERTIA\_GATE} & $N_{ig}$ & — & 0--20 & 0（关） & 纯 AC KLU 路径每 $N$ 次分解用 MUMPS 复检惯性；默认关闭（效果算例相关，注释 :solve\_primal\_dual\_ipm），且依赖本检出处不可用的 MUMPS（:solve\_primal\_dual\_ipm）\\
\srcpath{HACDCPF\_OPF\_RESCALE\_PER\_ITER} & — & — & 置位/未置位 & 未置位 & 置 1 恢复逐迭代 Ruiz 重平衡（默认冻结）（:factor\_assembled\_kkt）\\
\srcpath{HACDCPF\_OPF\_FILTER\_ALL} & — & — & 置位/未置位 & 未置位 & 对任意规模强制启用残差滤波器（默认仅 $n+m_{eq}\ge512$）（:solve\_primal\_dual\_ipm）\\
\srcpath{HACDCPF\_OPF\_STRICT\_THETA} & — & — & 置位/未置位 & 未置位 & 诊断：仅按纯可行下降接受（:solve\_primal\_dual\_ipm、:eval\_and\_test\_trial（lambda））\\
\srcpath{HACDCPF\_OPF\_NO\_LEGACY} & — & — & 置位/未置位 & 未置位 & 诊断：关闭遗留三轴接受路径（:solve\_primal\_dual\_ipm、:eval\_and\_test\_trial（lambda））\\
\srcpath{HACDCPF\_OPF\_PROF} & — & — & 置位/未置位 & 未置位 & 进程级分解/回代/装配耗时剖析，退出时打印（:IpmProf）\\
\srcpath{HACDCPF\_DISABLE\_IPOPT\_OPF} & — & — & 置位/未置位 & 未置位 & 运行期强制禁用嵌入式 Ipopt（src/optimal\_power\_flow/ac\_opf.cpp:solve\_parity\_with\_ipopt）\\
\end{paramtable}

\subsubsection*{内部硬编码常数}
\begin{paramtable}
\fld{kDenseAutoKktDim} & — & — & — & 512 & $n+m_{eq}$ 不大于此值走稠密 PartialPivLU（:solve\_primal\_dual\_ipm）\\
\fld{kLargeKktFilterThreshold} & — & — & — & 512 & 启用残差滤波器的 KKT 维数下限（:solve\_primal\_dual\_ipm）\\
\fld{kMaxBacktracks} & — & — & — & 14 & 主线搜索最大回溯次数（:solve\_primal\_dual\_ipm）\\
\fld{kBacktrack} & — & — & — & 0.5 & 回溯折半因子（:solve\_primal\_dual\_ipm）\\
\fld{kFilterEta} & $\eta$ & — & — & $10^{-4}$ & 遗留三轴进步余量（:solve\_primal\_dual\_ipm、:eval\_and\_test\_trial（lambda））\\
\fld{kGammaTheta}/\fld{kGammaPhi} & $\gamma_\theta$/$\gamma_\varphi$ & — & — & $10^{-5}$ & 滤波器占优余量（:solve\_primal\_dual\_ipm）\\
\fld{kDelta}\allowbreak /\fld{kSTheta}\allowbreak /\fld{kEtaPhi}\allowbreak /\fld{kThetaMin} & $\delta_s$/$s_\theta$/$\eta_\varphi$/$\theta_{\min}$ & — & — & 1.0/\allowbreak 1.1/\allowbreak $10^{-8}$/\allowbreak $10^{-4}$ & switching 与 Armijo 常数（:solve\_primal\_dual\_ipm）\\
\fld{kRegLevels} & $\delta$ & — & — & 见说明 & 凝集式渐进正则化梯级 $\{0,\,10^{-8},\,10^{-6},\,10^{-4}\}\lVert W\rVert$（:solve\_primal\_dual\_ipm）\\
— & $\delta_W$ & — & — & $10^{-8}\lVert L_{xx}\rVert$ 起、$\times8$、封顶 $10^{-2}\lVert L_{xx}\rVert$ & 增广式惯性校正（仅 MUMPS；本检出处不生效）（:solve\_primal\_dual\_ipm）\\
— & — & — & — & $10^{-12}$ & $z,\mu$ 下限保护（:solve\_primal\_dual\_ipm）与步长下限 $10^{-10}$（:solve\_primal\_dual\_ipm）\\
— & — & — & — & $10^{-7}$ & 空步判定阈值（:solve\_primal\_dual\_ipm）\\
\fld{kBetaGC} & $\beta_{gc}$ & — & — & 0.1 & Gondzio 额外校正目标 $0.1\gamma$（:solve\_primal\_dual\_ipm）\\
\fld{kStagnationWindow} & — & — & — & 20 & 目标停滞早停窗口（:solve\_primal\_dual\_ipm）\\
— & — & — & — & $100\times$ & 可接受级收敛的容忍度倍率（:solve\_primal\_dual\_ipm）\\
— & — & — & — & 5 & Ruiz 平衡最大轮数、$10^{-3}$ 停则（:compute\_ruiz\_scaling）\\
— & — & — & — & 2 & KKT 迭代精化步数（:kkt\_solve\_scaled、:kkt\_solve\_dense）；Gondzio 校正次数上限（:solve\_primal\_dual\_ipm）\\
— & — & — & — & 100 & 目标梯度缩放参考范数 $\lVert\nabla f\rVert_\infty$（:solve\_primal\_dual\_ipm）\\
\end{paramtable}

\subsection{验证与校核}\label{subsec:pipm-verification}

\begin{itemize}
  \item \textbf{单元/回归测试}：\srcpath{tests/test\_opf\_solver\_backends.cpp}
        以 \fld{ACOPFSolverBackend::ParityIPM} 覆盖后端分派、
        \fld{solver\_path} 标记、热启动状态回灌、不可行诊断与
        \fld{profiling} 计数；另覆盖纯 AC case30 和 Hybrid microgrid 的
        primal/dual 证书、零分解预算、兼容状态绕过，以及 case300/case2000
        的单调有界非认证路径。case300 还断言未接受 Phase I 时 Phase II 的
        初始残差与关闭 Phase I 完全一致；
        \srcpath{tests/test\_multiscale\_comprehensive.cpp} 与
        \srcpath{tests/test\_nighttime\_opf.cpp} 亦引用本路径；
  \item \textbf{基准算例}：源码注释记录了 case2000\_acdc（约 2~s）、
        case2869pegase（Release 54 迭代）、case13659 等 PEGASE/大电网
        算例的实测轨迹（:kkt\_form\_preference、:backend\_order），手算 OPF 算例与跨引擎
        对照通道见 \ref{sec:verification}~节；
  \item \textbf{雅可比正确性}：Parity 建模层的解析导数经有限差分核对
        （见 \ref{sec:parityform}~节），本核的 KKT 右端与恢复公式与其
        共享同一组回调（:solve\_primal\_dual\_ipm、:eval\_and\_test\_trial（lambda））。
  \item \textbf{性能口径}：\srcpath{phase\_one\_restoration\_benchmark}
        固定两次预热、五次计时。默认参数下 case30 从
        $3.1061\times10^{-2}$ 恢复到 $4.0305\times10^{-4}$，Phase II 从
        17/16 次迭代/分解降到 16/15，端到端中位数
        $4.458\to4.225$~ms（$1.055\times$）。Hybrid microgrid 的初始违反量
        $0.144>\mu_0$，case300 为 $1.129>\mu_0$，二者均由 admission gate
        以零分解跳过；Phase II 的迭代、分解与最终残差和关闭 Phase I 完全一致，
        结构探针中位成本分别约 $0.008$ 与 $0.124$~ms；对应端到端开/关
        中位数为 $0.235/0.228$~ms 与 $230.058/230.462$~ms。降低 $\mu_0$ 到
        0.03/0.01 会使 case30 也在 admission 外，因此默认 0.1 保留。
\end{itemize}

\subsection{边界与局限}\label{subsec:pipm-limits}

\begin{enumerate}
  \item \textbf{MUMPS 分支为死代码}：\srcpath{HACDCPF\_HAVE\_MUMPS} 在本
        检出处无任何定义点，\srcpath{MumpsSolver} 类不存在于当前
        MIPSolvers 头文件树；增广式的 $\delta_W$ 惯性校正、惯性门复检
        与注释中的 MUMPS 性能数字对当前构建均不适用，混合 AC/DC 问题实际
        由 UMFPACK/KLU/Eigen SparseLU 承接，且在 LU 类后端下惯性校验被
        静默跳过（:solve\_primal\_dual\_ipm）。
  \item \textbf{选项字段未生效}：\fld{IPMOptions::tol\_complementarity}
        与 \fld{regularization} 从未被读取；\fld{tol\_dual} 不进收敛
        判据（判据统一用 $\varepsilon_p$）。调用方按字段语义调参不会
        改变核行为。
  \item \textbf{小系统无全局化}：$n+m_{eq}<512$ 且未置
        \srcpath{HACDCPF\_OPF\_FILTER\_ALL} 时，线搜索无条件接受首个
        有限试步（:eval\_and\_test\_trial（lambda）），发散只能由 \texttt{KKT factorization failed}
        等异常路径兜底。
  \item \textbf{互补容忍度宽松}：互补判据下限 $2\times10^{-3}$（:solve\_primal\_dual\_ipm），
        故“收敛”不保证紧互补；严格 KKT 点需依赖对偶最小二乘认证或
        acceptable 级判据的对偶分量。
  \item \textbf{单线程}：核为单线程实现，剖析器按此假设设计
        （:IpmProf 注释）。
  \item \textbf{凝集式条件数}：默认凝集形式的 $\kappa(W)$ 随障碍比平方
        增长，刚性混合算例可能出现“精确求解但方向错误”的停滞；增广式
        可缓解但需自行开启（:KktForm、:kkt\_form\_preference）。
  \item \textbf{Ipopt 路径的平台限制与口径差异}：嵌入式 Ipopt 的可用性取决于
        MIPSolvers 检出是否含可构建的内嵌 \texttt{ipopt/} 源码（见本章门控节，
        默认仅 Apple 开启）；其 acceptable 判据（$1000\times$）
        比本核（$100\times$）宽松，且不返回约束乘子，LMP 在该路径不可用
        （src/optimal\_power\_flow/ac\_opf.cpp:solve\_parity\_with\_ipopt）。
  \item \textbf{悬置声明}：\srcpath{parity\_ipm\_solver.hpp:solve\_ac\_opf\_parity} 的
        \srcpath{solve\_ac\_opf\_parity} 无定义，调用入口实为
        \srcpath{solve\_ac\_opf} 内的静态驱动（src/optimal\_power\_flow/ac\_opf.cpp:solve\_with\_parity\_ipm）。
  \item 恢复阶段与中心性恢复仅存在于增广形式；默认凝集形式在滤波器
        全拒时直接报 \texttt{filter line search failed}（:solve\_primal\_dual\_ipm）。
\end{enumerate}
